\chapter{Case Library}
\label{ch:cases}

\section{Why the Book Needs a Research-Disciplined Case Library}

Without cases, the framework remains too close to a persuasive essay. Without
evidence discipline, the cases risk becoming confident storytelling. This
chapter therefore occupies a deliberately intermediate layer. It does
\emph{not} claim that ordinary study failure, exercise avoidance, or bedtime
drift have already been measured with one-to-one neuroscientific precision.
What it does claim is narrower and more useful:
\begin{enumerate}[nosep]
  \item the variables of the framework provide a stable analysis language;
  \item the research spine from the empirical, emergence, and advanced
        dynamical chapters supplies disciplined mechanism analogies;
  \item application-layer interventions become stronger when they are tied to
        explicit state-transition structure rather than vague motivational talk.
\end{enumerate}

The chapter therefore uses a three-lane reading rule for every case:
\begin{itemize}[nosep]
  \item \textbf{Framework lane:} What are the dominant variables?
  \item \textbf{Mechanism lane:} What research-backed state, path dependence,
        metastability, or attractor ideas are relevant?
  \item \textbf{Evidence lane:} Which claims are strong, which are medium, and
        which would become speculative if pushed too far?
\end{itemize}

\section{Case 1: Study Start Failure}

\subsection*{Situation}
A student returns from lunch intending to read a difficult chapter but instead
opens short-video feeds for forty minutes.

\subsection*{Framework map}
\begin{itemize}[nosep]
  \item \textbf{Current state:} post-lunch passive novelty consumption.
  \item \textbf{Target state:} seated proof-focused reading with active note
        production.
  \item \textbf{Dominant variables:} high $\Reward$, high $\Switch$, and high
        $\Env$.
  \item \textbf{Missed decision window:} immediate re-entry interval after the
        meal, before the phone reconstructs the old regime.
\end{itemize}

\subsection*{Mechanism interpretation}
This case is best read as a local state-transition failure. The current state
has an easily re-entered attractor-like character: low effort, rapid novelty,
and a familiar posture-tool configuration. The target state is not absent only
because the student lacks values; it is weakly stabilized and poorly prepared.
The emergence chapter adds a second layer: repeated post-lunch phone use can
create path dependence, so the same temporal boundary starts to function as a
history-loaded cue rather than a neutral moment.

\subsection*{Evidence lane discipline}
The strong empirical consciousness sources support only a limited transfer
here: organized states and transitions can exhibit structured stability and
changing dynamics. The case therefore uses \textbf{strong-to-medium}
translation, not direct proof. The sleep-transition bifurcation paper is
relevant as a model for why transition timing matters, but it does not prove
that every lunch break is a bifurcation in the strict neuroscientific sense.
Metastability language is useful in a \textbf{medium} sense: the work state
must become stable enough to hold, yet still be easy to recover after
interruption.

\subsection*{Operational repair}
\begin{enumerate}[nosep]
  \item Write the exact next proof target before the meal break.
  \item Leave the notebook open at the unfinished line.
  \item Remove the phone from the immediate work surface.
  \item Define the first post-break action as ``sit, read the last two lines,
        and write one continuation sentence.''
\end{enumerate}

\subsection*{Transfer lesson}
Study-start failure is usually not solved by increasing moral pressure. It is
solved by making the desired re-entry trajectory shorter, more visible, and
less dependent on fresh design at the moment of return.

\section{Case 2: Exercise Non-Start After Work}

\subsection*{Situation}
An employee endorses regular exercise every morning but rarely reaches the gym
after a draining workday.

\subsection*{Framework map}
\begin{itemize}[nosep]
  \item \textbf{Current state:} end-of-work depletion plus home-directed drift.
  \item \textbf{Target state:} arrival at the gym and completion of a prepared
        session.
  \item \textbf{Dominant variables:} low $\Will$, elevated $\DV$, and
        environment-dependent re-locking into a rest regime.
  \item \textbf{Missed decision window:} commute boundary before home entry.
\end{itemize}

\subsection*{Mechanism interpretation}
This case illustrates how timing and path dependence compound one another. If
the commute reliably terminates in the home-rest sequence, the trajectory is
already partially selected before conscious deliberation begins. In the
language of advanced dynamics, the home route behaves like a basin-reentry
channel. In the language of emergence, the repeated work-to-home-to-collapse
sequence becomes a historically consolidated transition policy.

\subsection*{Evidence lane discipline}
The current research spine supports a \textbf{medium} interpretation that
stability architecture and transition timing matter. It does \emph{not}
support a strong claim that evening exercise failure has been mapped onto a
specific neural attractor with clean empirical identification. The biologically
grounded lesson is narrower: state persistence, stability, and transition cost
are serious system properties, so intervention should target those properties
first.

\subsection*{Operational repair}
\begin{enumerate}[nosep]
  \item Move the decision upstream by packing the gym bag before work.
  \item Choose one default workout rather than a menu of options.
  \item Route the commute past the gym instead of directly home.
  \item Define success as physical entry plus one pre-specified minimum block
        of exercise, not an idealized full workout.
\end{enumerate}

\subsection*{Transfer lesson}
This case teaches a general rule: if the current route already reconstructs
the old state before deliberation starts, redesign the route rather than asking
the depleted self to overcome the full barrier in place.

\section{Case 3: Writing Avoidance in Research Work}

\subsection*{Situation}
A graduate student can read papers and discuss ideas but repeatedly delays
writing drafts.

\subsection*{Framework map}
\begin{itemize}[nosep]
  \item \textbf{Current state:} browsing, annotating, and discussing without
        visible output.
  \item \textbf{Target state:} rough written production with explicit claims.
  \item \textbf{Dominant variables:} moderate $\Reward$ in pseudo-work, very
        high $\Switch$ due to ambiguity, and self-protective environmental
        drift toward low-exposure tasks.
  \item \textbf{Missed decision window:} first minute of the writing session,
        when the draft could still be defined narrowly.
\end{itemize}

\subsection*{Mechanism interpretation}
This is a neighbouring-attractor case rather than a simple inactivity case.
The student is active, but active in a basin that provides intellectual
legitimacy without forcing externalized structure. Emergence language matters
because repeated success in reading or discussing can produce a path-dependent
identity loop: the system becomes efficient at the pre-draft regime and
therefore increasingly expensive to shift into visible production. Attractor
language helps here because the problem is not zero movement; it is movement
within the wrong region of state space.

\subsection*{Evidence lane discipline}
The mechanism claim here is \textbf{medium}. The attractor and metastability
literature supports using basin, persistence, and reconfiguration language as a
formal scaffold. It does not warrant a strong empirical claim that one can read
off ``writing avoidance'' directly from brain-level measurements. The case
therefore remains a disciplined behavioural model supported by systems
concepts, not a completed laboratory result.

\subsection*{Operational repair}
\begin{enumerate}[nosep]
  \item Redefine the target as one ugly but visible unit: five bullet claims,
        one paragraph skeleton, or one figure caption.
  \item Open the session with a production artifact, not a reading artifact.
  \item Keep one unresolved question from the previous session visible.
  \item Separate evaluation from generation so that sentence production does
        not have to carry full quality judgment at entry.
\end{enumerate}

\subsection*{Transfer lesson}
Some failures arise because the system has become too competent at an adjacent
state. In such cases the intervention must reduce target ambiguity and expose
the smallest real output unit, rather than merely extending time spent near the
target domain.

\section{Case 4: Sleep Schedule Breakdown}

\subsection*{Situation}
A person wants to sleep earlier but continues consuming stimulating content at
night until bedtime shifts later and later.

\subsection*{Framework map}
\begin{itemize}[nosep]
  \item \textbf{Current state:} high-stimulation late-night scrolling or media
        consumption.
  \item \textbf{Target state:} a designed pre-sleep regime that reliably
        transitions into sleep.
  \item \textbf{Dominant variables:} high $\Reward$, low resistance to
        continuation, weakly specified target state, and an environment that
        keeps brightness and stimulation active.
  \item \textbf{Missed decision window:} the first shutdown cue before the
        reward stream has fully captured the night.
\end{itemize}

\subsection*{Mechanism interpretation}
This is the case with the strongest direct connection to the research spine.
The sleep-transition literature gives a real example of patterned state change
that can be described in dynamical language. The manuscript should still stay
disciplined: the cited sleep work supports the importance of transition
structure and control parameters, not a full behavioural theory of personal
bedtime discipline. Even so, the case strongly benefits from thinking in terms
of regime change: if bright screens, novelty, and open-ended content keep the
system in the wrong basin, then ``just stop'' is not a real control strategy.

\subsection*{Evidence lane discipline}
This case has the strongest \textbf{strong} anchor available in the current
source set, because the indexed Nature Neuroscience sleep paper directly treats
falling asleep as a patterned transition dynamic. The translation from that
result to practical bedtime routine design remains partly \textbf{medium},
because the manuscript is still moving from biological regime transition to
everyday behavioural control design. That move is defensible, but it should be
stated as interpretation rather than finished proof.

\subsection*{Operational repair}
\begin{enumerate}[nosep]
  \item Replace the negative rule ``stop consuming'' with a positive pre-sleep
        state.
  \item Specify the object set: dim light, paper book, charger away from bed,
        and one fixed shutdown sequence.
  \item Treat the first shutdown cue as the key decision window rather than the
        final exhausted minute.
  \item Reduce parameter drift by keeping time, light, and device placement
        stable across evenings.
\end{enumerate}

\subsection*{Transfer lesson}
The more a target state is defined as the absence of a bad state, the weaker
its entry dynamics become. Productive control usually requires a replacement
regime, not only a prohibition.

\section{Cross-Case Analytical Matrix}

The point of a case library is not to tell four stories. It is to create a
reusable classification surface.

\begin{center}
\begin{tabularx}{\textwidth}{@{}l l X X@{}}
\toprule
Case & Dominant variables & Path dependence / emergence & Type \\
\midrule
Study start &
high $\Reward$, high $\Switch$, high $\Env$ &
post-lunch cue repeatedly rebuilds distraction regime &
re-entry failure \\

Exercise non-start &
low $\Will$, high barrier, route-dependent $\Env$ &
commute repeatedly selects the home-rest sequence &
boundary failure \\

Writing avoidance &
high ambiguity, moderate pseudo-work reward &
adjacent low-exposure work becomes historically efficient &
neighbouring-basin failure \\

Sleep breakdown &
high night reward, weak target-state identity, strong $\Env$ &
repeated late-night stimulation stabilizes a delayed shutdown regime &
replacement-state failure \\
\bottomrule
\end{tabularx}
\end{center}

To preserve the analytical content without overloading the table, the
cross-case dynamics and repair logic are summarized immediately below.

\begin{itemize}[nosep]
  \item \textbf{Study start.} Dynamics language: rapid re-entry into a familiar
        basin plus weak work-state restabilization. Highest-value repair:
        preload the exact next task and redesign the post-break surface.
  \item \textbf{Exercise non-start.} Dynamics language: boundary-sensitive
        trajectory selection; home acts like a basin channel. Highest-value
        repair: move the decision upstream and reroute entry before home
        capture.
  \item \textbf{Writing avoidance.} Dynamics language: movement inside a
        neighbouring attractor rather than total inactivity. Highest-value
        repair: shrink the output unit and begin with visible production.
  \item \textbf{Sleep breakdown.} Dynamics language: regime persistence and
        transition-threshold sensitivity. Highest-value repair: install a
        positive shutdown state and protect the first cue.
\end{itemize}

\section{Case Typology}

These four cases are intended to form a small analytical library:
\begin{enumerate}[nosep]
  \item \textbf{Re-entry failure:} the target state was previously active, but
        re-entry cost is too high.
  \item \textbf{Boundary failure:} the system misses the only useful transition
        interval and then gets captured by the old route.
  \item \textbf{Neighbouring-basin failure:} effort exists, but it is trapped
        in an adjacent state that feels productive while blocking the actual
        target.
  \item \textbf{Replacement-state failure:} the agent tries to stop the old
        regime without constructing a stable new one.
\end{enumerate}

The benefit of this typology is transfer. A reader who understands the type no
longer needs to memorize the anecdote. The case becomes a reusable diagnostic
pattern.

\section{Recommended Reading}

\begin{readingbox}
\textbf{Foundation}
\begin{itemize}[nosep]
  \item John Holland, \emph{Emergence}, for how repeated local interaction can
        consolidate higher-level patterns.
  \item Eugene Izhikevich, \emph{Dynamical Systems in Neuroscience}, for
        attractor, basin, and transition language.
  \item Stanislas Dehaene, \emph{Consciousness and the Brain}, for a readable
        bridge from state-based neuroscience to practical interpretation.
\end{itemize}
\textbf{Modern Research}
\begin{itemize}[nosep]
  \item \emph{Metastability demystified: the foundational past, the pragmatic
        present and the promising future} (Nature Reviews Neuroscience, 2024,
        medium).
  \item \emph{Associative conditioning in gene regulatory network models
        increases integrative causal emergence} (Communications Biology, 2025,
        medium).
  \item \emph{Cortical connectivity, local dynamics and stability correlates of
        global conscious states} (Communications Biology, 2025, strong as
        biology background, medium for behavioural transfer).
  \item \emph{Falling asleep follows a predictable bifurcation dynamic}
        (Nature Neuroscience, 2025, strong for sleep-state transition).
\end{itemize}
\textbf{Reading rule}
\begin{itemize}[nosep]
  \item Use the modern sources to discipline the mechanism language of the
        cases.
  \item Use the classic textbooks to keep the cases teachable and connected to
        durable concepts rather than only recent paper terminology.
  \item Keep frontier consciousness-physics proposals outside the case layer;
        they are too speculative for ordinary behavioural diagnosis.
\end{itemize}
\end{readingbox}

\begin{summarybox}
The case library now does more than illustrate the framework. It functions as
an application-layer bridge between the research spine and practical transfer.
Each case identifies dominant variables, names the relevant path-dependent or
state-transition mechanism, marks the evidence lane clearly, and extracts a
repair pattern that can generalize beyond the anecdote. That is what makes the
chapter feel like an analytical library rather than a set of motivational
stories.
\end{summarybox}
